\documentclass{beamer}
\usetheme{Madrid}
\usepackage{amsmath,amssymb,booktabs,xcolor}
\usepackage{listings}

\lstset{
  language=Python,
  basicstyle=\ttfamily\scriptsize,
  keywordstyle=\color{blue},
  commentstyle=\color{gray},
  stringstyle=\color{orange!70!black},
  showstringspaces=false,
  breaklines=true,
  frame=single,
  backgroundcolor=\color{gray!8},
}

\title{Practical Session 4: LLM Agents and RAG in Finance}
\subtitle{Hands-On: ReAct Tracing, Tool Design, and RAG Evaluation}
\author{Juan F.\ Imbet}
\institute{Paris Dauphine -- PSL University}
\date{Large Language Models in Finance}

\begin{document}

% ============================================================
% FRAME 1: Title
% ============================================================
\begin{frame}
\frametitle{}
\titlepage
\end{frame}

% ============================================================
% FRAME 2: Session overview
% ============================================================
\begin{frame}
\frametitle{Session Overview}

\textbf{Duration:} 1 hour

\medskip
\begin{enumerate}
  \item \textbf{Recap} (5 min) --- PRA loop, ReAct format, memory types, RAGAS metrics
  \item \textbf{Problem 1} (15 min) --- Trace a ReAct agent on a filing Q\&A task
  \item \textbf{Problem 2} (10 min) --- Design a tool schema for a financial calculator
  \item \textbf{Case Study} (20 min) --- Build and evaluate a minimal RAG pipeline over
        a mock earnings call corpus; answer three diagnostic questions
  \item \textbf{Discussion} (5 min) --- When would you add a second agent? When would
        you add a hook?
  \item \textbf{Solutions} (5 min) --- Walkthrough of Problems 1 and 2
\end{enumerate}

\medskip
\begin{block}{What you will practise}
  Tracing agent reasoning steps; writing Pydantic tool schemas; building and evaluating
  a RAG pipeline; thinking through governance checkpoints.
\end{block}
\end{frame}

% ============================================================
% FRAME 3: Recap --- PRA loop and ReAct
% ============================================================
\begin{frame}
\frametitle{Recap: PRA Loop and ReAct}

\textbf{PRA loop} $(\pi, \mu, \phi)$:
\begin{itemize}
  \item $\phi$: observation + memory $\to$ updated memory
  \item $\pi$: memory $\to$ distribution over actions (implemented by the LLM)
  \item $\mu$: action + environment $\to$ next observation
\end{itemize}

\medskip
\textbf{ReAct trajectory:}
\[
  \tau = (T_1, A_1, O_1,\; T_2, A_2, O_2,\; \ldots,\; T_n, \textsc{Finish})
\]
Each $T_i$ is a natural-language thought. Each $A_i$ is a tool call.
Each $O_i$ is the tool's return value.

\medskip
\textbf{Key principle:} actions ground reasoning in verifiable external data.
$\hat{f}(t) \neq f(t)$ in general; calling the API returns the exact $f(t)$.
\end{frame}

% ============================================================
% FRAME 4: Recap --- memory types and RAGAS
% ============================================================
\begin{frame}
\frametitle{Recap: Memory and RAG Evaluation}

\textbf{Memory types:}
\begin{center}
\begin{tabular}{lcc}
\toprule
Type & Size bound & Persistence \\
\midrule
In-context & Context window & Session only \\
Retrieval (vector DB) & Unbounded & Persistent \\
Long-term (structured DB) & Unbounded & Persistent \\
\bottomrule
\end{tabular}
\end{center}

\medskip
\textbf{RAGAS metrics:}
\begin{itemize}
  \item \textbf{Faithfulness}: claims in answer supported by retrieved context
  \item \textbf{Answer Relevance}: answer semantically addresses the query
  \item \textbf{Context Precision}: retrieved chunks are relevant
  \item \textbf{Context Recall}: ground-truth claims are covered by context
\end{itemize}

Finance priority: \textbf{Faithfulness} first (unfaithful = hallucinated claim = compliance risk).
Faithfulness $< 0.85$ $\to$ block automated publication.
\end{frame}

% ============================================================
% FRAME 5: Problem 1 --- trace a ReAct agent
% ============================================================
\begin{frame}
\frametitle{Problem 1: Trace a ReAct Agent (15 min)}

A compliance analyst asks an agent:

\begin{quote}
\textit{``What credit risks does Microsoft disclose in its most recent 10-K, and
how do they compare to the risks it disclosed in the 2022 10-K?''}
\end{quote}

\textbf{Available tools:}
\begin{itemize}
  \item \texttt{edgar\_search(ticker, form, limit)} --- returns accession numbers and dates
  \item \texttt{edgar\_section(accession, section)} --- returns text of a named section
  \item \texttt{summarise\_risks(text, focus)} --- returns structured risk list
  \item \texttt{compare\_lists(list\_a, list\_b)} --- returns added, removed, changed items
\end{itemize}

\medskip
\textbf{Your task:} write the full ReAct trajectory $(T_1, A_1, O_1, \ldots, \textsc{Finish})$.
\begin{itemize}
  \item Include at least 4 steps before the \textsc{Finish} action
  \item Each thought must justify \emph{why} you choose the next action
  \item Make up plausible (but fictional) observation values
\end{itemize}
\end{frame}

% ============================================================
% FRAME 6: Problem 2 --- tool schema design
% ============================================================
\begin{frame}[fragile]
\frametitle{Problem 2: Design a Financial Tool Schema (10 min)}

You are building an LLM agent for a fixed-income desk. Design a Pydantic schema
for the following tool:

\medskip
\textbf{Tool:} \texttt{compute\_bond\_price}

Computes the present value of a bond given its cash flows and yield to maturity.
Should support both annual and semi-annual coupon conventions.

\medskip
\textbf{Your schema must include:}
\begin{itemize}
  \item Face value (in USD, any positive value)
  \item Annual coupon rate (as a decimal)
  \item Years to maturity (positive integer)
  \item Yield to maturity (as a decimal; annual, not semi-annual)
  \item Coupon frequency (annual or semi-annual)
\end{itemize}

\medskip
\textbf{For each field:} write a \texttt{Field(description=...)} with units, valid range,
and at least one example value. Use \texttt{ge}, \texttt{le}, or \texttt{gt} constraints.

\medskip
\textit{Bonus: what is the formula for bond price? Add it as a docstring.}
\end{frame}

% ============================================================
% FRAME 7: Case study overview
% ============================================================
\begin{frame}
\frametitle{Case Study: RAG over Earnings Call Transcripts (20 min)}

\textbf{Setup:} You have a mini-corpus of 5 earnings call excerpts (provided as strings).
You will build a minimal RAG pipeline and evaluate it using RAGAS-style metrics.

\medskip
\textbf{The corpus} (fictional, for illustration):
\begin{itemize}
  \item \texttt{apple\_q4\_2023.txt} --- 3 paragraphs, CFO guidance + China exposure
  \item \texttt{microsoft\_q2\_2024.txt} --- 3 paragraphs, Azure growth + AI capex
  \item \texttt{jpmorgan\_q1\_2024.txt} --- 3 paragraphs, NIM guidance + credit losses
  \item \texttt{nvidia\_q3\_2023.txt} --- 3 paragraphs, data centre demand + export controls
  \item \texttt{berkshire\_2023.txt} --- 3 paragraphs, insurance float + equity portfolio
\end{itemize}

\medskip
\textbf{Queries to evaluate:}
\begin{enumerate}
  \item \textit{``What revenue guidance did Apple give for Q1 2024?''}
  \item \textit{``What did JPMorgan say about credit loss provisions?''}
  \item \textit{``Which company mentioned export controls?''} (cross-doc)
\end{enumerate}
\end{frame}

% ============================================================
% FRAME 8: Case study --- step 1, chunking and embedding
% ============================================================
\begin{frame}[fragile]
\frametitle{Case Study Step 1: Chunking and Embedding}

\begin{lstlisting}
from sentence_transformers import SentenceTransformer
import numpy as np

# --- Sample corpus (use provided strings in the lab notebook) ---
docs = {
    "apple_q4_2023": "Revenue guidance for Q1 2024 is $89-90B ...",
    # ... add remaining documents
}

# --- Chunk each document into sentences ---
def sentence_chunk(text):
    return [s.strip() for s in text.split('.') if len(s.strip()) > 20]

chunks = []
for doc_id, text in docs.items():
    for sent in sentence_chunk(text):
        chunks.append({"doc_id": doc_id, "text": sent})

# --- Embed all chunks ---
model = SentenceTransformer("all-MiniLM-L6-v2")
embeddings = model.encode([c["text"] for c in chunks])
\end{lstlisting}

\medskip
\textbf{Question A:} How many chunks does your corpus produce?
What is the mean chunk length in tokens?
\end{frame}

% ============================================================
% FRAME 9: Case study --- step 2, retrieval and generation
% ============================================================
\begin{frame}[fragile]
\frametitle{Case Study Step 2: Retrieve and Generate}

\begin{lstlisting}
from sklearn.metrics.pairwise import cosine_similarity

def retrieve(query, k=3):
    q_emb = model.encode([query])
    sims = cosine_similarity(q_emb, embeddings)[0]
    top_k = np.argsort(sims)[::-1][:k]
    return [chunks[i] for i in top_k]

def generate_answer(query, context_chunks, llm_fn):
    context = "\n".join(c["text"] for c in context_chunks)
    prompt = f"""Answer based ONLY on the context below.
If the answer is not in the context, say "Not found in context."

Context:
{context}

Question: {query}
Answer:"""
    return llm_fn(prompt)
\end{lstlisting}

\medskip
\textbf{Question B:} For query 1 (Apple guidance), which chunks are retrieved ($k=3$)?
Are they all from \texttt{apple\_q4\_2023}? What does this tell you about metadata filtering?
\end{frame}

% ============================================================
% FRAME 10: Case study --- step 3, RAGAS-style evaluation
% ============================================================
\begin{frame}
\frametitle{Case Study Step 3: Evaluate with RAGAS Metrics}

For each of the three queries, fill in this table:

\medskip
\begin{center}
\begin{tabular}{lcccc}
\toprule
\textbf{Query} & \textbf{Faith.} & \textbf{Ans.Rel.} & \textbf{Ctx.Prec.} & \textbf{Ctx.Rec.} \\
\midrule
Apple guidance & ? & ? & ? & ? \\
JPM credit losses & ? & ? & ? & ? \\
Export controls & ? & ? & ? & ? \\
\bottomrule
\end{tabular}
\end{center}

\medskip
Score each metric 0--1 manually using the definitions:
\begin{itemize}
  \item \textbf{Faithfulness}: does each sentence in the answer appear in the context?
  \item \textbf{Answer Relevance}: does the answer address the question?
  \item \textbf{Context Precision}: are the retrieved chunks actually useful?
  \item \textbf{Context Recall}: does the context contain everything needed to answer?
\end{itemize}

\medskip
\textbf{Question C:} Which query has the lowest faithfulness? Why?
What would you change in the pipeline to improve it?
\end{frame}

% ============================================================
% FRAME 11: Discussion
% ============================================================
\begin{frame}
\frametitle{Discussion: Agent Design Decisions (5 min)}

Discuss in pairs or small groups:

\medskip
\textbf{When would you add a second agent?}
\begin{itemize}
  \item When the task has two distinct sub-tasks with different tool requirements?
  \item When you need a critic or compliance checker separate from the generator?
  \item When latency can be reduced by parallelism?
\end{itemize}

\medskip
\textbf{When would you add a hook?}
\begin{itemize}
  \item When a cross-cutting concern (audit, cost monitoring) must apply to every action?
  \item When the concern is owned by a different team than the agent developer?
  \item When you need to enforce a policy without modifying the agent's core logic?
\end{itemize}

\medskip
\textbf{Governance question:}\\
Your earnings call agent produces a structured summary that a junior analyst will
use to write a client note. What HITL pattern would you impose, and why?
\end{frame}

% ============================================================
% FRAME 12: Solutions --- Problem 1
% ============================================================
\begin{frame}
\frametitle{Solutions: Problem 1 --- ReAct Trajectory}

\footnotesize
\begin{tabular}{p{0.1\textwidth}p{0.82\textwidth}}
\textbf{T1} & I need Microsoft's most recent 10-K and the 2022 10-K. I will search EDGAR for both. \\[3pt]
\textbf{A1} & \texttt{edgar\_search("MSFT", "10-K", 2)} \\[3pt]
\textbf{O1} & \texttt{[\{acc: "0001564590-24-...", year: 2024\}, \{acc: "0001564590-22-...", year: 2022\}]} \\[6pt]
\textbf{T2} & I have both accession numbers. Retrieve risk factors from each filing. \\[3pt]
\textbf{A2} & \texttt{edgar\_section("0001564590-24-...", "Risk Factors")} \\[3pt]
\textbf{O2} & \texttt{[~8,000 words of 2024 risk factors text]} \\[6pt]
\textbf{T3} & Now retrieve 2022 risk factors. \\[3pt]
\textbf{A3} & \texttt{edgar\_section("0001564590-22-...", "Risk Factors")} \\[3pt]
\textbf{O3} & \texttt{[~7,500 words of 2022 risk factors text]} \\[6pt]
\textbf{T4} & Summarise credit risks specifically from each set of risk factors. \\[3pt]
\textbf{A4} & \texttt{summarise\_risks(O2, "credit")} + \texttt{summarise\_risks(O3, "credit")} \\[3pt]
\textbf{O4} & Two structured risk lists. \\[6pt]
\textbf{T5} & Compare the two lists to identify what changed. \\[3pt]
\textbf{A5} & \texttt{compare\_lists(risk\_list\_2024, risk\_list\_2022)} \\[3pt]
\textbf{O5} & Added: \{AI model supply concentration\}; Removed: \{pandemic supply disruption\} \\[6pt]
\textbf{} & \textsc{Finish}(structured comparison with citations) \\
\end{tabular}
\normalsize
\end{frame}

% ============================================================
% FRAME 13: Solutions --- Problem 2
% ============================================================
\begin{frame}[fragile]
\frametitle{Solutions: Problem 2 --- Bond Price Tool Schema}

\begin{lstlisting}
from pydantic import BaseModel, Field
from typing import Literal

class BondPriceInput(BaseModel):
    """
    Bond price = sum_{t=1}^{n} C/(1+y)^t + F/(1+y)^n
    where C = coupon payment, y = period yield, n = total periods, F = face value.
    """
    face_value: float = Field(
        description="Face value in USD (e.g. 1000.0)", gt=0.0
    )
    annual_coupon_rate: float = Field(
        description="Annual coupon rate as a decimal (e.g. 0.05 for 5%)",
        ge=0.0, le=1.0
    )
    years_to_maturity: int = Field(
        description="Years to maturity (e.g. 10)", gt=0
    )
    yield_to_maturity: float = Field(
        description="Annual YTM as a decimal (e.g. 0.06 for 6%)",
        ge=0.0, le=1.0
    )
    frequency: Literal["annual", "semiannual"] = Field(
        description="Coupon frequency: 'annual' or 'semiannual'",
        default="semiannual"
    )
\end{lstlisting}
\end{frame}

% ============================================================
% FRAME 14: Key takeaways
% ============================================================
\begin{frame}
\frametitle{Key Takeaways}

\begin{enumerate}
  \item \textbf{ReAct grounds reasoning in evidence.}
        Each tool-call observation replaces a potentially hallucinated approximation.
        Thought steps explain \emph{why} each action is chosen.

  \item \textbf{Tool schemas are a first-class design artifact.}
        Field descriptions propagate into the model's context.
        Pydantic constraints reduce the error surface.

  \item \textbf{RAG evaluation must be decomposed.}
        Faithfulness (generation quality) and context recall (retrieval quality)
        fail for different reasons and require different fixes.

  \item \textbf{Metadata filtering is not optional.}
        Unconstrained semantic search is noisy for named-entity queries.
        Add document metadata at ingestion time.

  \item \textbf{Governance patterns match reversibility.}
        Pre-execution approval for irreversible actions.
        Post-hoc audit for high-volume, low-stakes tasks.
\end{enumerate}
\end{frame}

\end{document}
